%!TEX root = ../../main.tex
\section{기존 접근과 남은 어려움}
\label{sec:intro-prior}

MOBA 승률 연구는 경기의 승자를 목표로 삼아 왔다. 픽과 플레이어 이력 같은 경기 전 입력 \cite{costa2021,hitar2023}, 분 단위로 요약한 경기 중 상태 \cite{silva2018continuous}, 프로 경기의 실시간 예측 \cite{hodge2021win}, 리그 시즌의 그래프 모형 \cite{bisberg2022}이 그 예이며, 이런 예측기의 보정도 연구되었다 \cite{kim2020confidence}. 학습된 승률 모형은 킬과 아이템처럼 개별 사건의 가치를 그 사건이 더하는 승률로 매기는 통화로도 쓰였다 \cite{maymin2021smart,jalovaara2024win}. 이 계열은 경기 단위의 목표에 답하며, 교전 단위의 결과를 정의하는 데 쓸 수 있는 도구를 제공한다.

교전 단위 연구는 전투를 검출하고 그 발생이나 승자를 예측해 왔다. Dota~2의 encounter는 리플레이의 유닛 근접과 피해로 검출되었고 \cite{schubert2016encounter}, 공개 처리기는 영웅 사망을 중심으로 전투 창을 열며 \cite{opendota2018processteamfights}, 이후 연구는 사망이나 한타를 몇 초 앞서 예측하거나 \cite{tot2021camera,katona2019time} 검출한 전투를 경기 예측의 입력으로 썼다 \cite{ke2022}. 리그 오브 레전드의 공개 Match-V5 기록에서는 Baek과 Kwon \cite{baek2026killconditioned}이 킬 군집으로 교전을 정의하고 교전 전 관측으로 교전의 승자를 예측하였다. 이들 연구의 정의와 한계는 제 \ref{sec:rel-encounter} 절에서 다룬다.

이 과제는 세 가지 이유로 어렵다. 첫째, 공개 텔레메트리는 비동기이다. 상태는 1분 프레임으로, 사건은 밀리초로 기록되고 피해 사건은 없으므로 교전의 경계는 킬로부터 재구성해야 하며, 재구성한 경계에 대한 참 경계 주석은 없다. 둘째, 교전의 결과는 관측되지 않는다. 킬 수나 골드 교환은 관측되지만 교전이 경기 승패에 미친 영향은 관측되지 않으므로, 경기와 연결된 결과를 쓰려면 예측 전에 고정된 상태 평가기가 필요하고, 그 결과는 평가기에 조건부인 측정이 된다. 셋째, 현재 승률과 경기 시각이 이미 많은 것을 설명한다. 교전 전 상태의 예측 이득은 그 두 양의 유연한 함수를 넘어서 재야 하고, 경기 안에서 교전이 군집을 이루므로 경기 단위의 짝 비교로 검정해야 하며, 다른 패치와 지역에서 유지되는지 확인해야 한다 \cite{ChristoffersenDiebold2006,GneitingRaftery2007,CameronMiller2015}.

본 논문은 이 세 어려움에 대해 교전 정의의 상수를 자료와 게임 규칙에서 가져오고, 경기 승패로 학습해 동결한 평가기로 교전 구간의 승률 변화 방향을 결과로 정의하며, 현재 승률과 시간의 유연한 기준선을 구성하여 동일 교전에서 사전 상태의 추가 예측 효용을 평가한다.
